\section{Control de exportaciones y geopolítica}

\subsection{La evolución de las restricciones a los chips}

El control estadounidense de exportaciones de chips de IA hacia China atravesó varias etapas. Al principio se adoptó un criterio de dos factores (cómputo + ancho de banda de interconexión), luego se simplificó a una restricción pura de cómputo, y más recientemente se amplió a las más generales «reglas de difusión de IA».

\begin{table}[H]
\centering
\begin{tabular}{lp{8cm}}
\toprule
\textbf{Chip} & \textbf{Características} \\
\midrule
H100 & Especificaciones completas: no se permite exportar a China \\
H800 & Mismo cómputo que el H100, pero con el \textbf{ancho de banda de interconexión} recortado: DeepSeek optimizó en torno a esto \\
H20 & Cómputo reducido, pero con una relación de ancho de banda/capacidad de memoria mejor que la del H100: resulta más adecuado para la inferencia \\
\bottomrule
\end{tabular}
\caption{Comparación de los chips de NVIDIA sujetos a control de exportaciones}
\end{table}

En 2024, NVIDIA envió 1 millón de unidades del H20 a China, pero en 2025 canceló 2 millones de pedidos (anticipando una prohibición total inminente).

Una ironía clave: gracias a su ventaja en ancho de banda/capacidad de memoria, el H20 en realidad resulta \textbf{superior} al H100 en ciertos escenarios para modelos de inferencia (modelos tipo R1 u O1 que requieren un KV Cache considerable). El gobierno históricamente solo controló el cómputo (flops), pero la memoria y la interconexión son igual de importantes.

El control de exportaciones evolucionó desde el criterio inicial de dos factores (cómputo + ancho de banda de interconexión) hacia una restricción pura de cómputo, y recientemente se amplió a las más generales «reglas de difusión de IA» — que limitan a las entidades vinculadas a China a rentar menos de $<$2,000 GPU o comprar menos de $<$1,500 GPU. Pero cada ronda de restricciones dio origen a nuevas estrategias de evasión e innovaciones técnicas.

Desde la perspectiva de la cadena de suministro de semiconductores, el control de exportaciones creó una situación singular: NVIDIA se vio obligada a diseñar chips «capados» (H800, H20) para el mercado chino, y esos chips capados terminaron siendo el punto de partida de las optimizaciones de DeepSeek. Restringir el ancho de banda de interconexión $\to$ DeepSeek desarrolló optimizaciones de comunicación de bajo nivel en CUDA $\to$ esas optimizaciones también sirven para el H100, de modo que la «optimización extrema sobre hardware restringido» puede resultar más eficiente que la «práctica estándar sobre hardware ilimitado». Esta es una consecuencia imprevista del control de exportaciones.

\subsection{Argumentos y contraargumentos sobre el control de exportaciones}

El argumento de Dario Amodei: si la IA llega a ser extraordinariamente poderosa, quien la construya primero tendrá la ventaja militar; las democracias deberían ir a la delantera.

\begin{importantbox}{La distinción clave entre entrenamiento e inferencia}
\textit{``Entrenar un modelo en sí mismo casi no logra nada — desplegar el modelo para generar crecimiento económico, capacidad militar... eso sí requiere una enorme cantidad de cómputo.''} (Dylan)

El control de exportaciones limita principalmente la capacidad de China para \textbf{desplegar IA a gran escala}, no para entrenar un solo modelo. DeepSeek demostró justamente eso: aunque el hardware de entrenamiento estaba restringido, lograron entrenar un modelo de vanguardia. Pero el servicio de inferencia a gran escala (para atender a cientos de millones de usuarios) requiere una cantidad de GPU decenas de veces mayor que el entrenamiento.
\end{importantbox}

La línea de tiempo de la AGI según Nathan: después de 2030. Dylan, en cambio, considera que las capacidades podrían llegar entre 2027 y 2028, pero que el costo del despliegue será prohibitivo.

\subsection{La línea de tiempo de la AGI y la Guerra Fría de la IA}

Los directores ejecutivos de IA dicen ``2 años'' — sumando algunos años más, eso equivale aproximadamente a 2030. Pero las restricciones físicas limitan la fantasía del ``despliegue de AGI con un solo clic'': no se puede cambiar el mundo al día siguiente de terminar de entrenar un modelo, porque el despliegue, la integración y la adaptación requieren tiempo.

El episodio de DeepSeek podría marcar el inicio de una verdadera Guerra Fría de la IA. Tras la reunión de Liang Wenfeng (梁文锋) con el segundo líder más importante de China, el país anunció un plan de subsidios a la IA de un billón de yuanes (unos \$160B) — una respuesta más de tres veces superior a los \$50B de la CHIPS Act estadounidense.

Dylan ofrece un marco que invita a reflexionar sobre el efecto a largo plazo del control de exportaciones:

\begin{warningbox}{La paradoja temporal del control de exportaciones}
\textit{``En un mundo con un crecimiento económico del 2-3\%, el control de exportaciones garantiza que China gane a largo plazo. El control de exportaciones solo tiene sentido si la IA provoca una transformación enorme en el corto plazo.''} (Dylan)

La cadena lógica: si la AGI no llega en 5 años, el control de exportaciones solo retrasa a China unos años, y al final China la alcanza y se vuelve más autónoma. Si la AGI llega en 3 años, el control de exportaciones sí limita la capacidad de despliegue de China, y entonces cobra verdadero sentido. Pero el control de exportaciones dio origen a optimizaciones extremas al estilo DeepSeek — China podría lograr el mismo efecto con menos hardware.

Un riesgo más profundo: alejar a China de la tecnología de vanguardia aumenta el incentivo para emprender una acción militar contra Taiwán. El control de exportaciones apunta hacia economías futuras separadas — una vez formadas, son sumamente difíciles de revertir. La capacidad industrial de China en centros de datos y electricidad \textbf{supera ampliamente} a la de Estados Unidos. Una regularidad histórica: la paz se asocia con la hegemonía global unipolar; un orden multipolar equivale a inestabilidad.
\end{warningbox}

\subsection{Sistemas militares y autónomos}

La experiencia real de los drones en Ucrania muestra que la operación humana sigue siendo muy superior a los sistemas completamente autónomos. El avance tecnológico no significa que las armas autónomas sean viables de inmediato — el caos del entorno de batalla supera con creces al de un laboratorio.

La ciberguerra probablemente llegará antes que los enjambres de robots físicos. La ruta más probable para que la AGI se use como arma no es un ejército de robots como en las películas de ciencia ficción, sino la ingeniería social (operaciones de influencia de precisión impulsadas por IA), la desinformación (generación masiva de contenido indistinguible de lo real) y los ataques a infraestructura crítica (la red eléctrica, los sistemas financieros, las redes de comunicación).

\textit{``Un apagón en todo Estados Unidos durante dos días... eso provocaría asesinatos y caos.''} (Lex) Esto es una amenaza más cercana y más real que cualquier robot autónomo.

\subsection{Resumen de la sección}

El efecto del control de exportaciones es complejo: restringió el hardware de entrenamiento, pero dio origen a arquitecturas más eficientes; el H20 resultó ser mejor para la inferencia; y una dinámica de guerra fría está tomando forma. La contradicción central: limitar la capacidad de despliegue de China frente a impulsar su innovación autónoma frente a aumentar el riesgo militar en el estrecho de Taiwán.

